\clearpage\thispagestyle{empty}
\pdfbookmark[0]{Quellen zu Pallas (4. November 1813)}{pallas1813nov-sources}
\begin{center}\textbf{Quellen zu Pallas (4. November 1813)}\end{center}
\noindent Diese KI-Übertragung ins heutige Deutsch folgt dem Bericht mit der Überschrift Pallas in den \href{https://gdz.sub.uni-goettingen.de/id/PPN236007467}{\textit{Werken}, Band~VI} von Gauß (Göttingen, 1874), S.~369--370. Er erschien zuerst unter der Überschrift Göttingen in den \textit{Göttingischen gelehrten Anzeigen}, Stück~176 (4. November 1813), S.~\href{https://api.digitale-sammlungen.de/iiif/image/v2/bsb10538684_00729/full/1600,/0/default.jpg}{1753}, \href{https://api.digitale-sammlungen.de/iiif/image/v2/bsb10538684_00730/full/1600,/0/default.jpg}{1754}, \href{https://api.digitale-sammlungen.de/iiif/image/v2/bsb10538684_00731/full/1600,/0/default.jpg}{1755}. Die Links führen zum digitalisierten Exemplar der Bayerischen Staatsbibliothek.
\par\smallskip\noindent Der Punkt nach „Opposition“ im ersten Absatz und die Verweise in eckigen Klammern folgen der Werkausgabe. Der Erstdruck kürzt in der Vergleichstabelle wiederholte Beobachternamen ab und verwendet Dezimalkommas. In beiden Ausgaben steht vor der ersten beobachteten Deklination kein Vorzeichen. Im gesonderten Brief vom 31. August steht dort ein Pluszeichen. KI-Übertragung und KI-Satz.
\clearpage
